
\begin{abstract}
This tutored project explores the use of data mining and machine learning techniques to analyze and predict housing market dynamics in the United States.
The main objective is to study the influence of macroeconomic variables such as inflation, unemployment, and interest rates on housing prices, and to build predictive models capable of estimating future market trends.

To achieve this goal, a complete data science workflow was implemented following the CRISP-DM methodology.
The project includes data understanding, cleaning, preprocessing, exploratory analysis, predictive modeling, evaluation, and dashboard deployment.
Two machine learning models were selected for experimentation: Ridge Regression and Random Forest.
Their performance was assessed using standard evaluation metrics including Mean Absolute Error, Root Mean Squared Error, and the coefficient of determination.

In addition, an interactive dashboard was developed with Streamlit in order to visualize the dataset, compare periods, inspect correlations, and produce predictions.
The results show that machine learning models can capture important economic relationships and support the interpretation of housing market movements.
Random Forest provided the strongest predictive accuracy, while Ridge Regression remained useful as an interpretable baseline model.

This project demonstrates the value of combining data mining, machine learning, and interactive visualization in economic analysis and decision support.

\keywordss{Data Mining, Machine Learning, Housing Market, Prediction, Dashboard, Streamlit, Random Forest, Ridge Regression}
\end{abstract}

\bigskip
\textbf{Integrated project scope.}
The project is designed as a broader housing intelligence platform. In addition to exploratory analysis, regression modeling, and forecasting components, the system integrates OLAP-style multidimensional analysis, an AI-assisted chatbot, prompt-based analytical interaction, unsupervised learning, reinforcement-learning simulation, experiment tracking, model registry concepts, production-readiness checks, and paper-review comparison with real housing market theory. The Streamlit application therefore becomes not only a visualization tool but also an interactive decision-support environment that connects data preparation, machine learning, explainability, deployment, and business interpretation.

\textbf{Complementary keywords:} OLAP, AI chatbot, prompt engineering, unsupervised learning, reinforcement learning, experiment tracking, model registry, production readiness, paper review, Streamlit deployment.
